\chapter{Model checking with TLA+}
\label{ch:model_checking}

Some behaviour depends on the order in which concurrent operations occur. Tests exercise
selected runs of the venue, but cannot cover every possible order of thread steps, message
arrivals, failures, and restarts. The project also describes selected concurrent algorithms
in TLA+ and checks those descriptions with TLC, the TLA+ model checker.

\section{What model checking checks}

A TLA+ specification describes the state of a system, the actions that can change that
state, and properties that must hold. TLC starts from the specified initial state and
explores every state reachable under the chosen configuration. A configuration sets finite
bounds, such as the number of threads, messages, failures, or leadership changes.

A property that must hold in every reachable state is a safety property. Examples include
never using a slab after it has been destroyed and never allowing two instances to act as
leader at the same time. A liveness property describes progress, such as eventually
choosing a leader. Liveness checks depend on assumptions about which actions continue to
occur and how messages are delivered; the assumptions are part of the model and its
configuration.

When TLC finds a state that violates a property, it reports the sequence of actions that
reaches that state. This sequence is a counterexample. A counterexample makes the
failure concrete: it shows which operations can occur, in which order, and where the
property stops holding.

\section{Models in the project}

\subsection{High availability}

The high-availability models examine leadership and replication behaviour. The current
leadership model, \texttt{MajorityLeaseHA.tla}, describes two instances and a third voter.
It checks properties including whether two instances can act as leader at once and whether
they can lead under the same epoch. It also checks progress under stated assumptions about
message timing, fairness, and restarts.

\texttt{FollowerBehindHA.tla} checks the rule that an instance must not take the lead when
its matching engine has acted on commands that the instance does not hold. This model
treats leadership as an abstract decision and concentrates on the commands in each log,
the statements that restrict a follower from leading, and the events that write or copy
those records. The high-availability section of this book describes the properties,
results, and limits of these models in more detail.

\subsection{Framework algorithms}

The framework models focus on algorithms shared by multiple venue components:

\begin{itemize}
  \item \texttt{SlabReclaim.tla} checks when an empty slab can be queued and destroyed,
  including whether any thread can still access it.
  \item \texttt{QueueOfEmptySlabs.tla} checks the linked queue used to return empty slabs
  for reclamation, including whether a slab can be lost or taken more than once.
  \item \texttt{PoolFreeList.tla} checks the concurrent free list used by the fixed-size
  memory pool, including the counter that prevents an ABA error and the validation of
  free requests at the allocator boundary.
  \item \texttt{FollowerHandOver.tla} checks that a follower catching up with a
  write-ahead log receives every released record in order when responsibility passes
  from the reading thread to the thread that owns the log.
\end{itemize}

The detailed specifications, configurations, properties, and results are recorded in
\texttt{docs/availability/tla/findings.md} and
\texttt{docs/framework/tla/findings.md}.

\section{Checking that a model can find faults}

A run in which TLC finds no violation is useful only if the model and its properties are
capable of exposing a fault. The project therefore includes configurations that remove or
change a rule and expect TLC to find the resulting violation. For example, the free-list
model is also checked without its counter; TLC must then find a sequence in which a slot
appears to be available while another thread still holds it. The slab queue and follower
handover models also have configurations that remove rules needed to preserve their
properties.

The build can reproduce the recorded counterexamples for the high-availability lease
models and generate a page that presents them as step-by-step traces. When Java and the
TLA+ tools are available and this build feature is enabled, installation fails if the
counterexamples no longer reproduce their expected violations. The separate
\texttt{scripts/check\_pool\_free\_list\_tla.py} command checks two free-list
configurations that must pass and one configuration that must find the expected
violation.

\section{What the results establish}

TLC checks the specification, not the compiled C++ implementation. A result applies to
the behaviour described by the model, for the bounds and assumptions in the configuration.
It does not establish that the specification matches every relevant part of the
implementation. The specification must be compared with the code and kept aligned with
it.

The framework models also abstract from details that TLC does not check. In particular,
they do not verify that the C++ memory ordering of each atomic operation is sufficient,
or that the compiler emits the instruction the model assumes. A model can check all
interleavings of its abstract actions without checking the hardware and language rules
that implement those actions.

An exhaustive result is therefore evidence about a defined model, not a guarantee about
every possible execution of the venue. The finite bounds, timing assumptions, and
behaviour left out of each model determine what its result can support. Those limits are
recorded alongside the individual specifications.
